%% ==========================================================================
%%  1  Introduction
%% ==========================================================================
\section{Introduction}
\label{sec:intro}

Discrete fair division studies the allocation of resources that cannot be
fractionally assigned among agents with individual preferences, and sits at the
interface of mathematical economics and computer science~\cite{BCE16,End17}.
The classical fairness criterion is \emph{envy-freeness}~\cite{Fol67,Var74}: no
agent should prefer the bundle assigned to another agent over her own. For
indivisible items an envy-free allocation need not exist, as the standard
one-item two-agent instance shows, and a large body of work therefore studies
relaxations such as envy-freeness up to one good (EF1)~\cite{Bud11,LMMS04}.

A second and complementary line of work retains \emph{exact} envy-freeness and
buys off the residual envy with money. Each agent $i$ receives, in addition to
her bundle, a subsidy $\subsidy_i \ge 0$ supplied by a third party, and one asks
for an allocation together with a subsidy vector under which every agent weakly
prefers her own bundle-plus-subsidy to that of every other agent. The question
is how much money this requires. Maskin~\cite{Mas87} answered it for
unit-demand agents with $m = n$: under the scaling in which no agent values any
single item above $1$, a total subsidy of $n-1$ suffices. Halpern and
Shah~\cite{HS19} moved to the multi-demand setting with arbitrary $m$, gave the
characterisation of \emph{envy-freeable} allocations that the whole line now
runs on, showed that a subsidy of $(n-1)m$ is necessary and sufficient in the
worst case when the allocation is handed to us, and conjectured that $n-1$
suffices when we may choose the allocation. Brustle et
al.~\cite{BDNSV20} proved that conjecture in the stronger per-agent form --- one
dollar to each agent suffices for additive valuations, with the allocation
simultaneously EF1 and balanced --- and showed that for general monotone
valuations $2(n-1)$ per agent suffices, so that the total subsidy is
$O(n^2)$ and, remarkably, \emph{independent of the number of items}. Barman
et al.~\cite{BKNS22} then improved the monotone bound by a factor of $n$ on the
dichotomous subclass: when every marginal
$v_i(S \cup \set{g}) - v_i(S)$ lies in $\set{0,1}$, there is an allocation whose
accompanying subsidy is exactly $0$ or $1$ for each agent, hence at most $n-1$
in total, computable in polynomial time in the value-oracle model. Their bound
assumes neither additivity nor submodularity nor even subadditivity, and it is
tight: with $n$ agents and a single unit-valued good, $n-1$ agents must each be
paid $1$.

\paragraph{Chores.}
Every result above is about \emph{goods}. This paper asks the mirrored question,
in which every item is a \emph{chore}: an object that an agent would rather not
hold, so that acquiring it weakly decreases her utility. Chores are not an
exotic variant. Shift rotations, on-call duty, teaching and refereeing load,
committee service, maintenance tasks and the liabilities attached to an estate
are all allocation problems in which the items are burdens rather than prizes,
and in all of them money is already the standard instrument of repair --- as
overtime pay, hardship allowance, or a course-release budget. If anything the
subsidy model is more natural for chores than for goods, since compensating
somebody for taking on unwanted work needs less justification than paying
somebody for the privilege of not receiving a prize.

Concretely, we study \emph{negative dichotomous} valuations: for every agent
$i$, every $S \subseteq \items$ and every $g \in \items \setminus S$,
\[
  v_i(S) - v_i(S \cup \set{g}) \;\in\; \set{0,1}.
\]
Equivalently, each marginal $v_i(S \cup \set{g}) - v_i(S)$ is $0$ or $-1$: every
item is a chore for every agent, and taking on one more chore costs an agent
either nothing or exactly one unit. This is the exact mirror of the model
of~\cite{BKNS22}, and we impose no further structure --- in particular the
valuations need not be additive, submodular, or subadditive.

Binary marginals are the natural model whenever a task either lands on an agent
as a fresh unit of burden or is absorbed by a commitment she has already made.
An engineer already staffing an on-call window absorbs a second incident in that
window at no extra cost; a referee already reading for a venue absorbs one more
paper from the same batch; a driver already routed through a district absorbs an
extra stop along the way. Whether a given chore is free depends in an arbitrary
way on the rest of the agent's assignment, which is exactly why the model must
tolerate non-additive and even non-subadditive valuations, and it is the
counterpart of the scheduling example used by~\cite{BKNS22} to motivate
dichotomous valuations for goods.

\paragraph{Why this is not a change of sign.}
It is tempting to expect that a chores result follows from the corresponding
goods result by negating the valuations. It does not, for three reasons that we
record here because they shape everything that follows.

First, money is one-signed. The subsidy vector is constrained to
$\subsidy \in \R^n_{+}$, and negating an instance negates the valuations but
leaves the subsidies where they are. The polyhedron of feasible subsidy vectors
is therefore a genuinely different object in the two models, and no formal
duality carries a bound across.

Second, the obvious transformation is available only for two agents. Writing
$\cost_i := -v_i$ for the cost form of a negative dichotomous valuation, the
complement $\tilde{v}_i(S) := \cost_i(\items) - \cost_i(\items \setminus S)$ is
again dichotomous in the sense of~\cite{BKNS22}, so the class is closed under
complementation. But the complement transformation respects the allocation only
when $\items \setminus A_j = A_i$, that is, only when $n = 2$; see
Remark~\ref{rem:complement}. For $n \ge 3$ there is no such reduction, and that is
where the content of the problem lies.

Third, and most tellingly, envy flows the other way. Under goods, envy
concentrates \emph{into} the agent holding the valuable items: many agents envy
one, and the tight instance of~\cite{BKNS22} pays $n-1$ agents one dollar each
so that they stop envying the single agent who received the good. Under chores,
envy emanates \emph{out of} the agent carrying the load: one agent envies many.
Since the minimum subsidy to an agent is the weight of the heaviest directed
path leaving her in the envy graph~\cite[Theorem~2]{HS19}, the two situations
are not interchangeable, and one should not assume in advance that the same
bound is the right target.

\paragraph{What is already known, and what is open.}
Two facts delimit the problem before any work is done, and we state them
explicitly so that the contribution can be measured against them.

On the one hand, existence of a bounded subsidy is not in question. Negative
dichotomous valuations are doubly monotone --- every item is a chore for every
agent, which is a special case of every item being a good or a chore for each
agent --- and they satisfy the two-sided normalisation
$\abs{v_i(S \cup \set{e}) - v_i(S)} \le 1$. The reduction of Kawase et
al.~\cite{KMSTY24}, which converts any EF1 allocation into an envy-free
one, therefore applies off the shelf and yields a subsidy of at most $n-1$ per
agent and $n(n-1)/2$ in total.
\verify{this containment and the resulting bound against Reading\_9 directly;
confirm in particular which EF1 convention~\cite{KMSTY24} requires, since the
goods-only and two-sided forms of EF1 differ once chores are present}

On the other hand, the lower bound of the goods model survives intact.

\begin{example}[The $n-1$ lower bound transfers]
\label{ex:lowerbound}
Take $n$ agents and $m = n-1$ chores, with $\cost_i(S) = \abs{S}$ for every
agent $i$ --- identical, binary, additive costs, which are in particular
negative dichotomous. Every complete allocation leaves at least one agent
empty-handed. Under the balanced allocation, in which $n-1$ agents hold one
chore each and one agent holds nothing, each of the $n-1$ loaded agents envies
the empty agent by exactly $1$ and no other envy is present, so the minimum
subsidy assigns $1$ to each loaded agent and $0$ to the empty one, for a total of
$n-1$. No allocation does better: any complete allocation gives some agent a
bundle of size $k \ge 1$ while some other agent is empty, and that agent's
heaviest outgoing path already has weight $k$.
\end{example}

So the target statement is precisely the analogue of~\cite[Theorem~4]{BKNS22}:
that under negative dichotomous valuations there is an allocation admitting a
subsidy vector $\subsidy \in \set{0,1}^n$, hence a total subsidy of at most
$n-1$, computable in polynomial time with value-oracle access. Example~\ref{ex:lowerbound}
shows this would be tight. Between the $n-1$ per agent that~\cite{KMSTY24}
already gives and the $\set{0,1}$ per agent that we are asking for lies a factor
of $n$ in the total, exactly the gap that~\cite{BKNS22} closed on the goods side.
Our objective is to close it on the chores side or to exhibit an instance showing
that it cannot be closed.

\paragraph{Our Results.}
We conjecture that the $\set{0,1}$ guarantee survives
(Conjecture~\ref{conj:main}) and prove it in three restrictions: binary additive
costs, where an explicit $O(nm)$ construction gives $\optsubsidy \in \set{0,1}^n$
and total subsidy at most $n-1$ (Theorem~\ref{thm:binadd}); identical costs,
where greedily loading a cheapest bundle suffices
(Theorem~\ref{thm:identical}); and $n = 2$, which reduces
to~\cite{BKNS22} by Remark~\ref{rem:complement}. None of these needs additivity
or submodularity beyond what the restriction itself imposes.

Our main results are negative and concern the two strategies a reader would try
first. The item-by-item insertion template of~\cite{BKNS22} cannot be completed:
Theorem~\ref{thm:obstruction} gives a three-agent instance with a legal partial
state and an unallocated chore such that every choice of receiving agent forces a
subsidy of $2$, although the full instance is solvable with no subsidy at all. A
correct proof must therefore be free to re-open already-allocated bundles. Nor
can one recover by selecting a welfare-maximising allocation instead:
Proposition~\ref{prop:msw-false} exhibits an instance whose \emph{unique}
utilitarian-optimal allocation needs a subsidy of $2$ while an allocation of
strictly higher total cost needs none, so a correct algorithm must also be free
to give up total cost. Both instances are small, explicit and machine-checked.

We also isolate three structural facts used throughout: an arc-update lemma
(Lemma~\ref{lem:arcupdate}) that reduces the entire goods/chore asymmetry to two
lines; a two-tier characterisation of $\set{0,1}$ solutions
(Proposition~\ref{prop:twotier}); and a size-shift transform
(Theorem~\ref{thm:sizeshift}) exhibiting the chore problem as exactly the goods
problem of~\cite{BKNS22} plus a cardinality-balancing requirement, which is the
precise sense in which Conjecture~\ref{conj:main} does not follow from that work
as a black box.

\paragraph{Related Work.}
The subsidy line for goods is described above. Within it, the dichotomous
subclass has attracted separate attention: Halpern and Shah~\cite{HS19} settle
binary additive valuations with a total subsidy of $n-1$, Goko et
al.~\cite{GIK21} obtain the analogous bound for binary submodular valuations
together with a strategy-proof mechanism, and Barman et al.~\cite{BKNS22} subsume
both by dropping every structural assumption beyond binary marginals. Dichotomous
preferences are themselves a well-studied restriction in social
choice~\cite{BMS05,BEF21}. On the computational side, Caragiannis and
Ioannidis~\cite{CI21} show that minimising the subsidy is NP-hard and give a
fully polynomial-time approximation scheme for a constant number of agents, and
Narayan et al.~\cite{NSV21} study the welfare cost of achieving envy-freeness
with transfers.

The paper closest to the present setting is that of Kawase et
al.~\cite{KMSTY24}, which is the only work in this line that admits chores at
all. Working with doubly monotone valuations under the two-sided normalisation,
it decouples the two jobs that~\cite{BDNSV20} performed together: rather than
constructing a specific allocation and showing that it is cheap to subsidise, it
takes \emph{any} EF1 allocation as input and bounds the subsidy needed to
convert it, obtaining $n-1$ per agent and $n(n-1)/2$ in total. Since EF1
allocations are known to exist for doubly monotone valuations, the wider class
comes along at no cost. That reduction is what makes the present question
well-posed rather than open-ended, and it is the baseline we are trying to beat.

The complementary repair --- no money, weaker fairness notion --- is pursued for
binary valuations by Bu et al.~\cite{BSY23}, who show that EFX allocations always
exist and can be computed in polynomial time for general binary valuations, thus
extending the matroid-rank result of Babaioff et al.~\cite{BEF21}. Under binary
valuations one therefore has a clean dichotomy on the goods side: either pay
$0$ or $1$ per agent and obtain exact envy-freeness~\cite{BKNS22}, or pay
nothing and obtain EFX~\cite{BSY23}. Whether either branch survives the passage
to chores is open; this paper addresses the first.

\todo{a literature sweep of the chores-side fair division literature is still
outstanding. Nothing in the corpus $R1$--$R9$ covers indivisible chores except
\cite{KMSTY24}, and the related-work paragraph above therefore says nothing
about EF1/EFX for chores, about the goods-and-chores literature, or about prior
uses of subsidies for task assignment. Add those readings to
\texttt{References/}, extend \texttt{paper\_map\_R1\_to\_R9.md}, and rewrite this
paragraph before circulating.}
